\section{Method}
\label{sec:method}

\paragraph{Method: \ours{}.} \ours{} \Review[Abhijit (AI Assisted)]{is a latent world model whose output is a}{`Method:' repeats the section title. drop the paragraph heading. \rid{ibjr}} \emph{motion} of the feature grids it has
observed.
%
For every future step and every patch, it predicts where in a transport window of observed features the patch's content
comes from, whether the patch stays or moves, and a correction for content that motion cannot explain, as the lower row
of \cref{fig:method} shows on a held-out DROID window.
%
\Cref{sec:setting,sec:arch} define the task and the predictor, \cref{sec:transport,sec:training} the transport output
and its loss, and \cref{sec:theory} shows why this output suits scenes whose content mostly moves.

\begin{figure}[t]
  \centering
  \includegraphics[width=\linewidth]{architecture.pdf}
  \caption{\textbf{\ours{} architecture.} Top: a predictor turns observed features and actions into one query per patch and step; \ours{} decodes all steps in parallel from observed features. Bottom: on a held-out window, each forecast is kept + moved + corrected observed features.}
  \label{fig:method}
\end{figure}

\subsection{Problem setting}
\label{sec:setting}
A frozen encoder $E$ maps an image to a grid of $N=G\times G$ patch features $\Z_t\in\R^{N\times C}$. Given $H$
observed grids $\Z_{-H+1:0}$, the actions between them $\mathbf{a}_{-H+1:-1}$, and a future action sequence
$\mathbf{a}_{0:K-1}$, a world model predicts $\hat\Z_{1:K}$. The forecast $\hat\Z_k$ may depend only on
$\mathbf{a}_{<k}$ (causality).
%
Features are standardised per channel with train-split statistics shared across positions.

\subsection{Architecture}
\label{sec:arch}
\textbf{Memory encoder.} Each observed patch feature is projected to width $d$ and receives a spatial embedding, a frame
embedding and an embedding of the action applied after that frame.
%
$L_e$ transformer layers over all $H\cdot N$ tokens give memory states $\mathbf{M}\in\R^{H\times N\times d}$.
%
\textbf{Action-prefix queries.} A unidirectional GRU over embedded actions gives a causal prefix state, to which a
horizon embedding is added to form $\mathbf{c}_k$.
%
Query tokens start from the memory states of the last frame, and $L_d$ decoder layers apply self-attention across
patches and cross-attention to $\mathbf{M}$, each modulated by $\mathbf{c}_k$ through adaptive layer normalisation (AdaLN) with
zero-initialised gains \citep{peebles2023dit}.
%
The result is a query $\h_{k,i}\in\R^d$ for patch $i$ at step $k$, and all $K$ steps are decoded in parallel.

\subsection{Gated transport}
\label{sec:transport}
A regression head would write every future feature from $\h_{k,i}$, whereas the transport reads it from the observed
grids.
%
For query patch $i$, let $\mathcal{W}(i)$ be the $w\times w$ spatial window around $i$ in each of the last $S\le H$
observed frames. Transport weights are
\begin{equation}
  \ReviewEquationBox[Abhijit (AI Assisted)]{\pi_{k,i}(j) \;=\; \operatorname{softmax}_{j\in\mathcal{W}(i)}\Big(
  \tfrac{1}{\sqrt{d_{\text{key}}}}\,(\mathbf{W}_{\text{qry}}\h_{k,i})^{\!\top}(\mathbf{W}_{\text{key}}\mathbf{M}_{j}) + \beta\,\mathbf{1}[j=(0,i)]\Big),
  \label{eq:transport}}{Abhijit (AI Assisted)}{simplify: attention over the window plus an identity bias $\beta$. say that in words first. \rid{rgj7}}
\end{equation}
where $\mathbf{W}_{\text{qry}},\mathbf{W}_{\text{key}}$ project to key width $d_{\text{key}}$, $j$ indexes a source patch and
frame, and $\beta$ is a learned identity bias toward the same location in the last frame.
%
The transport moves \emph{observed} features, not latent states, $\mathbf{T}_{k,i}=\sum_{j}\pi_{k,i}(j)\,\Z_j$.
%
A per-patch gate
$g_{k,i}=\sigma(\mathbf{w}_g^\top\h_{k,i}+b_g)$, with sigmoid $\sigma$, and a correction $\mathbf{r}_{k,i}=\mathbf{s}\odot\mathbf{W}_o\h_{k,i}$ give
\begin{equation}
  \ReviewEquationBox[Abhijit (AI Assisted)]{\hat\Z_{k,i} \;=\; (1-g_{k,i})\,\Z_{0,i} \;+\; g_{k,i}\,\mathbf{T}_{k,i} \;+\; \mathbf{r}_{k,i}.
  \label{eq:output}}{Abhijit (AI Assisted)}{keep. add that Direct is this with $g{=}0$, so the reader sees the one change. \rid{ymi7}}
\end{equation}
The first term keeps the last observed feature in place, the second moves observed content, and the third adds what
no observed feature contains.
%
$\mathbf{W}_o$ is zero-initialised and $b_g=-3$, so an untrained model starts close to persistence, and $\mathbf{s}\in\R^C$ is a learned per-channel scale.

\subsection{Training}
\label{sec:training}
We minimise the mean squared error over all horizons, $\mathcal{L}_{\text{pred}}=\frac1K\sum_k\|\hat\Z_k-\Z_k\|^2$.
%
Anchoring every forecast to observed features makes it easy to ignore the actions, since a forecast that stays close
to $\Z_0$ scores well on average.
%
An optional action-contrastive hinge with margin $m$ and stop-gradient $\operatorname{sg}$ counters this by comparing the error under the true future actions with the error
under actions borrowed from another sample in the batch,
\begin{equation}
  \ReviewEquationBox[Abhijit (AI Assisted)]{\mathcal{L}_{\text{act}} = \max\!\big(0,\; m - (\|\hat\Z(\tilde{\mathbf a})-\Z\|^2 - \operatorname{sg}\|\hat\Z(\mathbf a)-\Z\|^2)\big),
  \qquad \mathcal{L}=\mathcal{L}_{\text{pred}}+\lambda\,\mathcal{L}_{\text{act}} .
  \label{eq:contrastive}}{Abhijit (AI Assisted)}{$\lambda{=}0$ in every result and the ablation shows no gain. move this loss to the appendix. \rid{m528}}
\end{equation}
All forecasting and plug-in results use $\lambda{=}0$, like every baseline, so $\mathcal{L}_{\text{act}}$ enters only as
an ablation.

\subsection{Motion model}
\label{sec:theory}
\begin{wrapfigure}[19]{r}{0.47\linewidth}
  \vspace{-1.2em}
  \centering
  \includegraphics[width=\linewidth]{featurespace.pdf}
  \caption{\textbf{One patch in feature space} under the three forecast routes \ding{172}--\ding{174}.}
  \label{fig:featurespace}
  \vspace{-0.6em}
\end{wrapfigure}
We motivate \cref{eq:output} with a simple generative model, drawn for one patch in \cref{fig:featurespace}.
%
\ours{} mixes observed vectors (\ding{172}, Prop.~\ref{prop:bayes}), an additive head synthesises the displacement
(\ding{173}, Prop.~\ref{prop:residual}), and recursion feeds errors forward (\ding{174}, Prop.~\ref{prop:compound}).
%
Suppose the grid at step $k$ arises from motion plus new content,
\begin{equation}
  \ReviewEquationBox[Abhijit (AI Assisted)]{\Z_k \;=\; \mathbf{P}_k\,\Z_0 + \bm{\varepsilon}_k,
  \label{eq:motion-model}}{Abhijit (AI Assisted)}{simplify in words: each future patch copies one observed patch, plus zero-mean new content. \rid{qgxd}}
\end{equation}
where $\mathbf{P}_k\in\{0,1\}^{N\times N}$ is a row-stochastic \emph{motion} matrix (each patch $i$ copies one source
patch $\rho(i)$ at integer-patch offsets; $\mathbf{P}_k=\mathbf{I}$ is a static scene; sub-patch motion is covered by soft weights between neighbouring sources) and $\bm{\varepsilon}_k$ is new content with
$\E[\bm{\varepsilon}_k\mid\mathbf{P}_k,\Z_0,\mathbf{a}]=\mathbf{0}$. Motion is uncertain given the context $c=(\Z_{\le 0},\mathbf{a}_{<k})$.

\begin{proposition}[The optimal forecast is a transport]\label{prop:bayes}
Under \cref{eq:motion-model}, the MSE-optimal forecast is
$\E[\Z_k\mid c]=\bar{\mathbf{P}}_k\Z_0$ with $\bar{\mathbf{P}}_k=\E[\mathbf{P}_k\mid c]$, a row-stochastic matrix.
If the support of each row of $\mathbf{P}_k$ lies in the transport window $\mathcal{W}$, then $\bar{\mathbf{P}}_k\Z_0$
can be approximated to arbitrary precision by the transport of \cref{eq:transport} with $g{=}1$ and
$\mathbf{r}{=}\mathbf{0}$, given logits that the decoder can express.
\end{proposition}

A sufficiently expressive head could reach this optimum as well, but transport parameterises it exactly.
%
Real features also change through viewpoint, lighting and deformation, which the correction $\mathbf{r}$ absorbs.
%
An additive head $\Z_0+\mathbf{W}_o\h$ instead has to write out the displacement in feature space.

\begin{proposition}[Size of the regression target]\label{prop:residual}
For a deterministic motion $\mathbf{P}$ with sources $\rho(i)$, the residual an additive head must output is
$\mathbf{R}=(\mathbf{P}-\mathbf{I})\Z_0$ with
$\|\mathbf{R}\|_F^2=\sum_i\|\mathbf{z}_{0,\rho(i)}-\mathbf{z}_{0,i}\|^2$, the feature difference along the motion. The
residual required by \ours{} is $\mathbf{0}$ for every such $\mathbf{P}$ within the window.
\end{proposition}

The proposition concerns what each head must represent, not what it can learn, as an additive head can learn to copy
but its target grows with the motion and this feature difference.
%
On static patches ($\mathbf{P}=\mathbf{I}$), the identity path is exact, whereas a regression head must output exactly
zero, and any error corrupts the static scene.
%
Finally, \ours{} decodes every horizon from observed features, which avoids the compounding of recursive decoding;
\cref{app:proofs} gives all proofs.

\begin{proposition}[Recursive versus direct decoding]\label{prop:compound}
Let a recursive predictor $f$ be $L$-Lipschitz in its input with one-step error at most $\delta$. Its $k$-step
rollout error is bounded by $\delta\sum_{j=0}^{k-1}L^j$, which grows geometrically for $L>1$. A direct predictor with
per-horizon error $\delta_k$ \Review[Abhijit (AI Assisted)]{incurs exactly}{true by definition, and $\delta_k$ can grow with $k$. compare bounds on the same quantity or call this a remark. \rid{5r4r}} $\delta_k$, independent of its errors at other horizons.
\end{proposition}

